\clearpage
\chapter*{2026 Febrero}
\addcontentsline{toc}{chapter}{2026 Febrero}

\begin{enunciado}
    \hypertarget{2026 B 1}{}
    \textbf{1. (3 puntos)} Demostrar que si en un espacio metrico \(X\) toda sucesión infinita tiene un punto de acumulación, entonces el espacio \(X\) es compacto.
\end{enunciado}

\begin{enunciado}
    \hypertarget{2026 B 2}{}
    \textbf{2. (3 puntos)}\\
    - Consideremos el conjunto \(\mathbb{R}\) de los números reales con la topología usual \(T_u\). Definimos la relación de equivalencia \(W\) en \(\mathbb{R}\) como:
    \[
        xWy \Leftrightarrow x,y \in \mathbb{Q} \text{ ó } x = y
    \]
    donde \(\mathbb{Q}\) es el conjunto de los números racionales. Sea \(X = R/W\) el espacio cociente con la topología cociente.\\
    - Consideremos el espacio topológico producto \((\mathbb{R} \times \mathbb{R}, T_u \times T_u)\). Definimos la relación de equivalencia \(S\) en \(\mathbb{R} \times \mathbb{R}\) como:
    \[
        (x,y)S(x',y') \Leftrightarrow \{(x,y),(x',y')\} \subset \mathbb{Q} \times \mathbb{Q} \text{ ó } (x,y) = (x',y').
    \]
    Sea \(Y = (\mathbb{R} \times \mathbb{R})/S\) el espacio cociente con la topología cociente.\\
    - Estudiar si el espacio cociente \(Y = (\mathbb{R} \times \mathbb{R})/S\) es homeomorfo al espacio producto de cocientes \(X \times X = (\mathbb{R}/W) \times (\mathbb{R}/W)\).
\end{enunciado}

\begin{enunciado}
    \hypertarget{2026 B 3}{}
    \textbf{3. (4 puntos)} Consideremos en el conjunto \(\mathbb{R}\) de los números reales la topología \(T\) definida mediante \(T = \{A \subset \mathbb{R} \mid 0 \notin A \text{ o } (-1, 1) \subset A\}\). Estudiar
    \begin{enumerate}[label=\alph*)]
        \item Si el espacio \((\mathbb{R}, T)\) es \(T_2\).
        \item Si el espacio \((\mathbb{R}, T)\) es compacto.
        \item Si el espacio \((\mathbb{R}, T)\) verifica el segundo axioma de numerabilidad.
    \end{enumerate}
\end{enunciado}

